\chapter{相关技术与理论基础}
\label{chap:technology}

本章介绍系统涉及的关键技术。首先说明大语言模型的基本原理和提示工程方法，然后分析检索增强生成的流程与适用边界，接着讨论文本向量化、混合检索、重排序和多轮对话中的关键问题，最后介绍前后端分离架构下的身份认证与权限控制。这些内容构成后续需求分析和系统设计的理论基础。

\section{大语言模型与提示工程}
\label{sec:llm}

大语言模型通常以 Transformer 的解码器结构为基础，通过在海量文本上预测下一个词元进行预训练。模型在生成时根据已有上下文逐个产生词元，并利用注意力机制捕捉长距离依赖。经过指令微调和对齐训练后，模型能够按照用户的自然语言指令完成问答、摘要、改写和代码生成等任务。对于问答系统而言，大模型的价值主要在于语言理解和组织能力：它可以把用户口语化的问题转换为规范的表述，也可以把检索到的多条证据整合成连贯的回答。

大模型的输出具有概率性，受到温度、Top-$p$、最大生成长度等参数影响。温度越低，输出越稳定、保守；温度越高，输出越发散，但事实性风险也随之增加。校园问答要求答案准确、口径统一，因此需要把温度调低，并通过系统提示明确角色、任务、证据范围和输出格式。提示工程在这里的作用不是“让模型更聪明”，而是把任务边界说清楚：哪些内容必须依据证据，哪些内容不能推测，证据不足时应该如何回复，引用编号应该放在什么位置。对于结构化输出，还可以要求模型按照固定的 JSON 字段返回答案、引用编号和置信度，便于后端解析和校验。

本课题不训练或微调大模型，原因有三点。其一，训练大模型需要大规模语料和算力，超出一名本科生毕业设计的条件；其二，校园资料规模相对有限，通过检索增强即可把校本知识注入回答过程，没有必要把知识固化到参数中；其三，使用成熟的大模型接口可以让研究重点集中在知识库构建、检索融合、引用溯源和评价方法上，符合任务书“不要求训练或微调大模型”的要求。为避免对单一模型服务产生依赖，系统在模型接口层做抽象，配置项中只保存服务地址、模型名称和密钥引用，更换模型时不需要修改业务代码。

\section{检索增强生成}
\label{sec:rag}

检索增强生成把生成过程拆成“索引—检索—生成”三个阶段。索引阶段把外部文档切分成片段，计算向量或建立倒排索引，写入向量数据库或搜索引擎；检索阶段根据用户问题召回相关片段，并可选择性地进行融合和重排；生成阶段把问题、候选证据和输出约束组织成提示词，交给大模型生成答案。与纯参数化模型相比，RAG 的优势在于知识可以随时更新、答案可以附带出处、错误可以定位到具体资料。其基本流程如图~\ref{fig:rag-flow} 所示。

\begin{figure}[H]
  \centering
  \begin{tikzpicture}[
    node distance=0.55cm and 0.55cm,
    proc/.style={draw, rounded corners=3pt, minimum width=2.15cm, minimum height=0.85cm, align=center, fill=gray!8, font=\zihao{5}},
    arr/.style={-{Latex[length=2mm]}, thick}
  ]
    \node[proc] (q) {用户问题};
    \node[proc, right=of q] (rw) {查询处理};
    \node[proc, right=of rw] (ret) {混合检索};
    \node[proc, right=of ret] (ev) {候选证据};
    \node[proc, below=1.0cm of ev] (gen) {受约束生成};
    \node[proc, left=of gen] (cit) {引用映射};
    \node[proc, left=of cit] (ans) {带引用回答};
    \draw[arr] (q) -- (rw);
    \draw[arr] (rw) -- (ret);
    \draw[arr] (ret) -- (ev);
    \draw[arr] (ev) -- (gen);
    \draw[arr] (gen) -- (cit);
    \draw[arr] (cit) -- (ans);
  \end{tikzpicture}
  \caption{RAG 问答的基本流程}
  \label{fig:rag-flow}
\end{figure}

RAG 并不能自动消除幻觉。检索不到相关证据、召回内容与问题无关、证据之间互相矛盾、模型没有遵守提示词约束，都会导致回答失败。因此，一个可用的 RAG 系统需要在每个环节设置可控点：检索阶段要有召回率和排序质量指标，生成阶段要有证据约束和拒答策略，输出阶段要有引用校验。Self-RAG 通过反思标记让模型判断是否需要检索、证据是否相关、生成内容是否有支持\upcite{asai2024}，其思想可以借鉴，但其训练成本较高。本课题采用轻量化的替代方案：在检索后设置相关性阈值，在提示词中要求“只依据证据回答”，在输出后检查引用编号是否落在候选证据范围内，任一环节不满足条件时转为拒答或提示用户补充信息。

从知识组织方式看，RAG 也可以与知识图谱结合。知识图谱擅长表达实体之间的结构化关系，适合处理“某学院有哪些专业”“某课程的先修课程是什么”这类关系型问题；RAG 擅长处理非结构化文本中的解释性内容\upcite{dou2026}。校园问答以规章制度和办事指南为主，答案往往是一段说明性文字，因此本课题以文本片段检索为主，暂不引入知识图谱，以控制系统复杂度。

\section{文本向量化与语义检索}
\label{sec:embedding}

文本向量化把自然语言映射到固定维度的实数向量，使语义相近的文本在向量空间中距离较近。DPR 采用双塔结构，问题和段落分别经过独立的编码器，训练目标是让问题与正确段落的向量相似度高于错误段落\upcite{karpukhin2020}。在推理阶段，所有段落向量可以预先计算并建立索引，用户提问时只需编码问题，再在向量库中做近邻搜索。与 BM25 相比，稠密检索不要求问题与段落共享相同的词语，因此对同义改写、口语化表达和跨句语义匹配更友好。

中文向量化模型的选择需要考虑语料领域和部署条件。C-Pack 提供了面向中文的评测基准、训练语料和 BGE 系列模型，覆盖语义相似度、检索、聚类等多种任务\upcite{xiao2024}。本课题优先选用参数量适中的 BGE 中文模型，在本地 CPU 上完成知识片段和问题的编码，以降低对外部服务的依赖。对于长度超过模型输入上限的文档，需要先分块再编码；分块过粗会引入无关内容，分块过细则会破坏上下文，因此分块策略需要结合文档结构和使用场景确定。

在向量检索中，常用余弦相似度衡量问题向量与片段向量的接近程度。设问题向量为 $\mathbf{q}$，片段向量为 $\mathbf{d}$，其余弦相似度定义为
\begin{equation}
  \cos(\mathbf{q},\mathbf{d})=\frac{\mathbf{q}\cdot\mathbf{d}}{\lVert\mathbf{q}\rVert\,\lVert\mathbf{d}\rVert}.
  \label{eq:cosine}
\end{equation}
当向量已经归一化时，余弦相似度等价于内积，计算量更小。实际系统中通常取相似度最高的 $k$ 个片段作为候选，$k$ 的取值需要在召回率和上下文长度之间权衡。$k$ 过小可能漏掉关键证据，$k$ 过大则会把噪声带入生成阶段，增加模型混淆和响应时延。

向量数据库负责保存片段向量、原文和元数据，并支持按相似度检索。常见的实现包括 Chroma、Milvus、FAISS 等。Chroma 部署简单、提供 Python 接口，适合本科课题的原型开发；Milvus 更适合大规模数据和高并发场景。本课题选用 Chroma 作为默认实现，同时在数据访问层保留接口抽象，便于后续替换。除向量外，每个片段还需要保存文档编号、标题路径、页码、分类、更新时间和权限标签等元数据，这些信息既用于结果展示，也用于检索过滤和引用溯源。

\section{混合检索与重排序}
\label{sec:hybrid}

稀疏检索和稠密检索的互补性，是混合检索的依据。BM25 基于词频和逆文档频率计算相关性，对专有名词、文件编号、条款序号等精确匹配非常有效，但无法理解同义表达；稠密检索能够捕捉语义相似性，却可能忽略关键的数字、名称和编号。把两条通道的候选合并，可以在不显著增加成本的前提下提高召回率。BM25 的常用形式为
\begin{equation}
  \mathrm{BM25}(q,d)=\sum_{t\in q}\mathrm{IDF}(t)\cdot
  \frac{f(t,d)\,(k_1+1)}{f(t,d)+k_1\left(1-b+b\cdot\frac{\lvert d\rvert}{\mathrm{avgdl}}\right)},
  \label{eq:bm25}
\end{equation}
其中 $f(t,d)$ 表示词项 $t$ 在文档 $d$ 中的词频，$\lvert d\rvert$ 为文档长度，$\mathrm{avgdl}$ 为平均文档长度，$k_1$ 和 $b$ 为调节参数。BM25 不需要训练，计算速度快，适合作为混合检索中的精确匹配通道。

两条通道的候选不能直接按分数相加，因为 BM25 分数和余弦相似度的量纲不同。常用的融合方法是倒数排序融合，它只利用候选在各通道中的排名，不依赖原始分数：
\begin{equation}
  \mathrm{RRF}(d)=\sum_{r\in R}\frac{1}{k+\mathrm{rank}_r(d)},
  \label{eq:rrf}
\end{equation}
其中 $R$ 为召回通道集合，$\mathrm{rank}_r(d)$ 为文档 $d$ 在通道 $r$ 中的排名，$k$ 为平滑常数，通常取 60。RRF 的优点是实现简单、对分数尺度不敏感，适合异构检索通道的融合。融合后的候选还需要经过一次重排序。交叉编码器把问题和候选片段拼接后一起输入模型，能够建模两者之间的细粒度交互，排序精度通常高于双塔模型，但计算成本也更高，因此只对顶部若干条候选使用。

图~\ref{fig:dual-cross} 对比了双塔检索与交叉重排的差异。双塔结构可以预先编码文档，适合大规模召回；交叉编码器必须在线计算问题与每个候选的匹配分数，适合对少量候选做精排。本课题采用“双路召回—RRF 融合—交叉编码器重排—Top-$k$ 截断”的链路，在召回率和时延之间取得平衡。

\begin{figure}[H]
  \centering
  \begin{tikzpicture}[
    box/.style={draw, rounded corners=2pt, minimum width=2.3cm, minimum height=0.8cm, align=center, fill=gray!8, font=\zihao{5}},
    enc/.style={draw, rounded corners=2pt, minimum width=2.3cm, minimum height=0.8cm, align=center, fill=gray!16, font=\zihao{5}},
    arr/.style={-{Latex[length=2mm]}, thick}
  ]
    % 左侧：双塔
    \node[box] (q1) at (0,0) {问题};
    \node[enc] (qe) at (0,-1.0) {问题编码器};
    \node[box] (d1) at (3.2,0) {候选片段};
    \node[enc] (de) at (3.2,-1.0) {片段编码器};
    \node[draw, circle, minimum size=0.75cm, font=\zihao{5}] (sim) at (1.6,-2.1) {相似度};
    \draw[arr] (q1) -- (qe);
    \draw[arr] (d1) -- (de);
    \draw[arr] (qe) -- (sim);
    \draw[arr] (de) -- (sim);
    \node[font=\zihao{5}] at (1.6,-2.95) {双塔检索};

    % 右侧：交叉编码器
    \node[box] (q2) at (6.2,0) {问题};
    \node[box] (d2) at (9.2,0) {候选片段};
    \node[enc, minimum width=3.0cm] (ce) at (7.7,-1.2) {交叉编码器};
    \node[draw, circle, minimum size=0.75cm, font=\zihao{5}] (score) at (7.7,-2.2) {得分};
    \draw[arr] (q2) -- (ce);
    \draw[arr] (d2) -- (ce);
    \draw[arr] (ce) -- (score);
    \node[font=\zihao{5}] at (7.7,-2.95) {交叉重排};
  \end{tikzpicture}
  \caption{双塔检索与交叉重排的结构对比}
  \label{fig:dual-cross}
\end{figure}

表~\ref{tab:retrieval-compare} 从匹配方式、优点、局限和适用场景四个方面比较了三种检索方案。可以看出，混合检索并不是简单地“二选一”，而是把稀疏通道的精确性和稠密通道的语义泛化能力结合起来，再通过重排序提升顶部结果的相关性。

\begin{table}[H]
  \centering
  \caption{三种检索方案的对比}
  \label{tab:retrieval-compare}
  \zihao{5}
  \renewcommand{\arraystretch}{1.25}
  \begin{tabular}{p{2.2cm}p{3.4cm}p{3.6cm}p{3.6cm}}
    \toprule
    对比维度 & BM25 稀疏检索 & 稠密向量检索 & 混合检索 \\
    \midrule
    匹配方式 & 词项匹配、逆文档频率加权 & 语义向量相似度 & 双路召回后融合 \\
    主要优点 & 对专有名词、编号、条款敏感；无需训练 & 能处理同义表达和口语化提问 & 兼顾精确匹配与语义泛化 \\
    主要局限 & 难以处理同义改写；受分词影响 & 对数字、编号不敏感；依赖嵌入模型 & 参数增多，融合与重排带来额外时延 \\
    适用场景 & 文件编号、条款查询 & 概念解释、流程咨询 & 校园综合问答 \\
    \bottomrule
  \end{tabular}
\end{table}

\section{多轮对话与问题改写}
\label{sec:multi-turn}

校园咨询往往不是一问一答。用户可能先问“国家助学金的申请条件是什么”，接着追问“那需要哪些材料”，再问“什么时候截止”。如果系统只把当前问题送入检索，第二、第三个问题缺少主语和限定条件，检索结果很容易偏离主题。多轮对话管理的任务，就是把当前问题与最近的会话历史结合起来，改写成可以独立检索的完整问题，同时保留原始问题，以便追溯用户真实意图\upcite{gan2026}。

常见做法是保留最近若干轮对话，把用户当前问题、上一轮问题和系统回答一起送入改写模型，要求模型补全指代、省略和隐含条件，输出一个不依赖上下文的自包含问题。改写后的查询用于检索，原始问题用于最终回答和日志记录。这样既能提高召回率，又不会因为改写错误而丢失原始语义。对于“那研究生呢”这类对比式追问，改写结果应保留前文的主题，只替换限定对象；对于“这个流程要多久”这类指代式追问，改写结果应把“这个流程”替换为上一轮的具体事项。

多轮对话还涉及历史记录的管理。系统需要保存会话、消息、引用和反馈，支持用户查看历史、继续提问和重新发起会话。为了避免上下文过长，通常只保留最近 $N$ 轮对话，并对早期内容做摘要。摘要可以压缩 token 消耗，但也可能丢失细节，因此本课题仅在会话轮次较多时启用摘要，并在引用展示中保留原始证据，保证用户仍能核对答案来源。

\section{Web 开发与权限控制}
\label{sec:web-security}

系统采用前后端分离的 B/S 架构。前端使用 Vue 3 和 Element Plus 构建页面，负责登录、问答、引用展开、历史记录和管理后台等交互；后端使用 Python 与 FastAPI 提供 RESTful 接口，负责身份认证、业务逻辑、文档处理、检索和生成；MySQL 保存用户、角色、文档、会话、消息、反馈和配置等关系数据；向量数据库保存知识片段索引；大模型、嵌入模型和重排序模型通过统一的模型服务层访问。前后端通过 JSON 交换数据，接口使用统一的响应结构和错误码，便于前端处理异常。

身份认证采用 JWT。用户登录成功后，后端校验密码哈希，签发包含用户编号、角色和过期时间的令牌；前端在后续请求的 Authorization 头中携带令牌；后端通过依赖注入解析令牌、校验签名和有效期，并把当前用户信息注入请求上下文\upcite{lu2025}。JWT 是无状态的，服务端不需要保存会话，适合前后端分离和水平扩展；但令牌一旦签发，在过期前无法主动失效，因此需要设置较短的有效期，并配合刷新令牌或黑名单机制处理注销和权限变更。

授权采用基于角色的访问控制（Role-Based Access Control，RBAC）。系统定义学生、教师和管理员等角色，为角色分配接口权限，为用户分配角色。后端在处理请求前检查当前角色是否具备所需权限，前端根据权限隐藏不可操作的菜单和按钮。权限控制不能只做在页面层，否则用户可以直接调用接口绕过限制；因此文档上传、删除、统计查看等敏感操作必须在后端再次校验\upcite{wei2025}。对于知识库检索，还需要在检索阶段按文档的权限标签过滤，避免低权限用户通过问答间接获取受限内容\upcite{liao2026}。

除认证授权外，系统还需要处理以下安全问题：密码使用 bcrypt 或 Argon2 等算法加盐哈希存储，不保存明文；大模型和向量库的密钥通过环境变量或配置文件注入，不写入代码仓库；上传文件校验扩展名、MIME 类型和大小，解析时限制页数和超时时间，防止恶意文件导致服务不可用；接口对输入长度、分页参数和检索条数做范围校验，防止注入和资源耗尽；日志对手机号、身份证号等敏感字段脱敏，避免二次泄露。这些措施将在第~\ref{chap:design} 章的安全设计中进一步展开。
